\chapter{Koszul--Petersson Pairing Compatibility}
\label{v4-arith:kp-compat}

An arithmetic Koszul pairing can be compared with a Petersson form
only after their domains have been related. The notation
\(\mathcal V^{\mathrm{prim}}\) denotes the Hochschild trace class
of the identity on the proposed global arithmetic Iwahori--Hecke
category of \(GL_2\). A state realization of this class, its dual
coefficient space, and its automorphic placement are additional
constructions. Chapter~\ref{v4-arith:closure} proves positivity
conditional on those maps and the pairing identity.

The Segal and prime Hilbert models provide explicit positive forms.
An identification with the arithmetic bar--cobar pairing requires
the injective placement, local pairing maps and completed adelic
extension of Hypothesis~\ref{v4-arith:kms-gns:hyp:arithmetic-pairing}.
Every arithmetic Koszul--Petersson comparison in this chapter
retains these hypotheses. The prime cyclic GNS construction gives
an actual local map; it does not supply the arithmetic isometry.

\section{Statement Recalled}
\label{v4-arith:kp:statement}

\subsection{Setup}
\label{v4-arith:kp:setup}

Write \(\mathcal V^{\mathrm{prim}}_{\mathrm{Heis}}\) for the chosen
local tensor model
\begin{equation}
\label{v4-arith:kp:eq:Vprim-adelic}
\mathcal{V}^{\mathrm{prim}}_{\mathrm{Heis}}\;=\;
\mathcal{H}^{(\infty)}\;\otimes\;{\bigotimes_{p}}^{\!\prime}\,\mathcal{H}^{(p)},
\end{equation}
with archimedean Heisenberg Fock space
\(\mathcal{H}^{(\infty)}=\mathbb{C}[\alpha_{-1},\alpha_{-2},\ldots]\cdot|0\rangle_{\infty}\)
and finite-place factors
\(\mathcal{H}^{(p)}=C_{c}(GL_{2}(\mathbb{Z}_{p})\backslash GL_{2}(\mathbb{Q}_{p})/GL_{2}(\mathbb{Z}_{p}))\)
under convolution. Normalize Haar measure by
\(\operatorname{vol}(GL_2(\mathbb Z_p))=1\), and use
\(|0\rangle_p=\mathbf1_{GL_2(\mathbb Z_p)}\) as reference vector.
The form \(\int f(g)\overline{h(g)}\,dg\) is positive on this
space and gives the reference vector norm one. The restricted
tensor product uses these vectors, so its product form is positive
by the finite-tensor argument of Chapter~\ref{v4-arith:closure}.

Any additional operators \(\psi_p\) must be specified on their
actual local carriers. Pointed \(E_0\)-data in
Definition~\ref{v4-arith:def:E0-local-factor} do not determine
Hecke operations or a Bost--Connes representation. Equation
\eqref{v4-arith:kp:eq:Vprim-adelic} defines a tensor model.
Its identification with a state realization of the arithmetic
Hochschild class remains part of the comparison hypothesis.

An additional arithmetic dg algebra comparison would identify the
Koszul-dual carrier with the Bost--Connes Q-lattice system
\begin{equation}
\bigl(\mathcal{V}^{\mathrm{prim}}_{\mathrm{Heis}}\bigr)^{!}
\;=\;\mathbb{A}^{\mathrm{BC}}
\end{equation}
as asked in Question~\ref{v4-arith:cor:koszul-dual-BC}.
The Bost--Connes presentation alone does not provide that comparison.

Let \(\tau^{\mathrm{BC}}\) be the algebraic
Bost--Connes anti-involution,
\[
\tau^{\mathrm{BC}}(\mu_{n})=\mu_{n}^{*},\qquad
\tau^{\mathrm{BC}}(e(\gamma))=e(-\gamma),
\]
the \(*\)-operation of the Bost--Connes \(C^{*}\)-algebra
(Bost--Connes~\S1.2). On the prime cyclic GNS space at \(\beta=1\),
for \(m=p^k\), the modular lift satisfies
\[
 J\pi(\mu_m)\Omega=m^{1/2}\pi(\mu_m^*)\Omega.
\]
This is Proposition~\ref{v4-arith:kms-gns:prop:J-equals-sigma-BC}
on its Hilbert--Schmidt carrier. In arithmetic pairing formulae,
\(\sigma_p^{\mathrm{BC}}\) denotes the separately specified
antilinear map compatible with that model under the comparison
hypothesis. An action on additional phase generators requires the
chosen full state and its representation.

\subsection{The Two Pairings}
\label{v4-arith:kp:two-pairings}

Under the specified arithmetic carrier and dual-coefficient
comparisons, the Koszul evaluation takes the bilinear form
\begin{equation}
\label{v4-arith:kp:eq:Koszul-pairing-def}
\langle\,\cdot\,,\,\cdot\,\rangle^{\mathrm{Koszul}}\colon
\mathcal{V}^{\mathrm{prim}}_{\mathrm{Heis}}\otimes\mathbb{A}^{\mathrm{BC}}\;\longrightarrow\;\mathbb{C}
\end{equation}
obtained by transporting the chosen bar--cobar evaluation.
An adjunction \(\Omega^{\mathrm{Ar}}\dashv B^{\mathrm{Ar}}\)
alone does not identify its evaluation target with \(\mathbb C\)
or its dual coefficient space with \(\mathbb A^{\mathrm{BC}}\).
Those maps belong to the hypotheses. A self-pairing further requires
a map from the conjugate state space to that dual coefficient space.

Let \(V\) denote the chosen state realization, with injective
placement \(\iota\) into \(L^2([GL_2])\). The transported Petersson
form is
\begin{equation}
\label{v4-arith:kp:eq:Pet-pairing-def}
\langle\,\cdot\,,\,\cdot\,\rangle^{\mathrm{Pet}}
\colon V\otimes\overline V\;\longrightarrow\;\mathbb{C},
\quad
\langle v,w\rangle^{\mathrm{Pet}}
\;=\;\int_{[GL_{2}]}(\iota v)(\gamma)\overline{(\iota w)(\gamma)}\,d\gamma.
\end{equation}
Here \(d\gamma\) is the chosen Tamagawa measure on
\(Z(\mathbb{A})\,GL_{2}(\mathbb{Q})\backslash GL_{2}(\mathbb{A})/K_{\infty}\).
An adelic Whittaker realization
must be constructed and compared with the arithmetic carrier before
it supplies such a placement.

\subsection{The Identity}
\label{v4-arith:kp:identity}

\begin{theorem}[Koszul--Petersson compatibility, domain-restricted]
\label{v4-arith:kp:thm:main}
Assume the arithmetic placement, local pairing maps and completed
adelic extension of Hypothesis~\ref{v4-arith:kms-gns:hyp:arithmetic-pairing}
on the stated normalized core.
\emph{(domain-restricted to the normalized
sph-fin cuspidal core; see
Remark~\ref{v4-arith:kp:rem:domain-restriction})}. For all
\(v,w\in\mathcal{V}^{\mathrm{prim}}_{\mathrm{sph\text{-}fin,cusp}}\),
\begin{equation}
\label{v4-arith:kp:eq:main}
\langle v,w\rangle^{\mathrm{Pet}}
\;=\;
\langle v,\sigma^{\mathrm{BC}}(w)\rangle^{\mathrm{Koszul}}.
\end{equation}
Equivalently, the adelic Petersson pairing equals the Koszul
pairing composed with \(\mathrm{id}\otimes\sigma^{\mathrm{BC}}\)
under the identification
\((\mathcal{V}^{\mathrm{prim}})^{!}=\mathbb{A}^{\mathrm{BC}}\), on the
normalized sph-fin cuspidal core.
\end{theorem}

\begin{remark}[domain restriction and open sub-items]
\label{v4-arith:kp:rem:domain-restriction}
Theorem~\ref{v4-arith:kp:thm:main} is conditional on the arithmetic
placement and pairing comparisons, including on the normalized
sph-fin cuspidal core. The explicit constructions below give weighted
occupation, prime Gibbs, spherical Whittaker and archimedean vacuum-line
models. They do not identify the full arithmetic bar carrier with an
automorphic representation. That identification, its compatibility
with the local actions and conjugations, and its completed adelic
extension are the remaining comparison maps in
Hypothesis~\ref{v4-arith:kms-gns:hyp:arithmetic-pairing}.
Ramified factors and the Eisenstein continuous spectrum require their
own carriers and comparison maps; the spherical calculation does not
supply them.
\end{remark}

\section{Pairing Domains and Normalizations}
\label{v4-arith:kp:obstruction-analysis-table}

The Koszul pairing is bilinear, while Petersson is conjugate-linear
in its second variable. Composition with the specified antilinear
map has the required sesquilinearity. The completed arithmetic
carrier, its local actions and the Haar normalization must be
preserved by the comparison maps. At almost every finite place
they must preserve the selected reference vector, so their tensor
products are defined on the restricted tensor domain.

The prime Gibbs state and its modular conjugation are constructed
in Chapter~\ref{v4-arith:kms-gns-full}. They are not an adelic
Plancherel identification. Neither the existence of spherical
vectors nor the equality of scalar Euler factors constructs the
arithmetic pairing maps. The remaining comparison is
Question~\ref{v4-arith:kms-gns:question:arithmetic-pairing}.

\section{Archimedean Identification}
\label{v4-arith:kp:arch}

\subsection{The Fock Inner Product}
\label{v4-arith:kp:arch:Fock}

The archimedean factor
\(\mathcal{H}^{(\infty)}=\mathbb{C}[\alpha_{-1},\alpha_{-2},\ldots]\cdot|0\rangle_{\infty}\)
is the Heisenberg vacuum Fock module. Its canonical Hermitian
form is the \emph{Segal inner product}: the unique
positive-definite sesquilinear form \(\langle\,\cdot\,,\,\cdot\,\rangle_{\mathrm{Segal}}\)
on \(\mathcal{H}^{(\infty)}\) satisfying
\begin{equation}
\label{v4-arith:kp:eq:Segal-normalisation}
\alpha_{-n}^{\dagger}\;=\;\alpha_{n},
\qquad
\langle |0\rangle_{\infty},|0\rangle_{\infty}\rangle_{\mathrm{Segal}}\;=\;1,
\qquad
\alpha_{n}|0\rangle_{\infty}\;=\;0\;\;\text{for}\;n\geq 1,
\end{equation}
together with the Heisenberg CCR
\([\alpha_{m},\alpha_{n}]=m\delta_{m+n,0}\) (Segal 1981
Commun.~Math.~Phys.~80; Frenkel--Ben-Zvi 2004 \S3.5; Kac 1998
\emph{Vertex Algebras for Beginners} \S2.5 Prop.~2.2).

\begin{proposition}[Segal form on monomials]
\label{v4-arith:kp:prop:Segal-monomials}
For monomials
\(v_{1}=\alpha_{-n_{1}}\cdots\alpha_{-n_{k}}|0\rangle_{\infty}\),
\(v_{2}=\alpha_{-m_{1}}\cdots\alpha_{-m_{l}}|0\rangle_{\infty}\)
with \(n_{1}\geq\cdots\geq n_{k}\geq 1\),
\(m_{1}\geq\cdots\geq m_{l}\geq 1\),
\begin{equation}
\label{v4-arith:kp:eq:Segal-monomials}
\langle v_{1},v_{2}\rangle_{\mathrm{Segal}}
\;=\;
\delta_{k,l}\cdot\sum_{\sigma\in S_{k}}\,
\prod_{i=1}^{k}n_{i}\,\delta_{n_{i},m_{\sigma(i)}}.
\end{equation}
Equivalently, Segal's form is the \(L^{2}\)-inner product of the
Wick-normal-ordered fields.
\end{proposition}

\begin{proof}
Wick's theorem applied to a product of creation/annihilation
operators equates
\(\langle v_{1},v_{2}\rangle_{\mathrm{Segal}}
=\langle|0\rangle_{\infty},\alpha_{n_{1}}\cdots\alpha_{n_{k}}
\alpha_{-m_{1}}\cdots\alpha_{-m_{l}}|0\rangle_{\infty}\rangle\).
Normal-ordering and using the CCR to collect the contractions
gives exactly
(\ref{v4-arith:kp:eq:Segal-monomials}). See Kac 1998 \S2.5
Prop.~2.2 or Frenkel--Ben-Zvi 2004 \S3.5.4.
\end{proof}

\subsection{Bilinear and Hermitian Fock Forms}
\label{v4-arith:kp:arch:Koszul-infty}
On the algebraic Heisenberg Fock space, let \(J_\infty\) conjugate
coefficients in the monomial basis. Define the bilinear form
\begin{equation}
\label{v4-arith:kp:eq:Koszul-infty}
 B_\infty(v,w)=H_{\mathrm{Segal}}(v,J_\infty w).
\end{equation}
Its Gram matrix on monomials is the real permutation sum of
Proposition~\ref{v4-arith:kp:prop:Segal-monomials}.
\begin{proposition}[The Fock conjugation identity]
\label{v4-arith:kp:prop:arch-koszul-is-Segal}
The form extends continuously to the Segal Hilbert completion, and
\begin{equation}
\label{v4-arith:kp:eq:arch-match}
 B_\infty(v,J_\infty w)=H_{\mathrm{Segal}}(v,w).
\end{equation}
\end{proposition}
\begin{proof}
The monomial Gram matrix is real, so coefficient conjugation is an
isometric involution. Hence \(|B_\infty(v,w)|\le\|v\|\|w\|\),
which gives the extension, and \(J_\infty^2=1\) gives the identity.
Identifying an arithmetic bar--cobar pairing with this form requires
maps on its two actual carriers, preserving the pairing and its
operator domains.
\end{proof}
\begin{corollary}[Archimedean conjugation type]
\label{v4-arith:kp:cor:arch-involution}
A bilinear arithmetic pairing identified with \(B_\infty\) becomes
the Segal Hermitian form after the transported coefficient
conjugation in its second slot. That conjugation is antilinear; it
cannot be replaced by the complex-linear identity on all vectors.
\end{corollary}
\begin{proof}
Apply \eqref{v4-arith:kp:eq:arch-match}; replacing \(w\) by \(iw\)
also directly distinguishes the two types.
\end{proof}

\subsection{An Explicit Archimedean Whittaker Line}
\label{v4-arith:kp:arch:Pet-infty}
Fix the function \(W_0(a)=|a|^{1/2}K_0(2\pi|a|)\) and the measure
\(da/|a|\) on \(\mathbb R^\times\). This specifies a spherical
Whittaker function in the spectral-parameter-zero model; a different
spectral parameter has a different Bessel order. The Hermitian form is
\begin{equation}
\label{v4-arith:kp:eq:Pet-arch-Whittaker}
 H_{\mathrm{Whit},\infty}(W,W')=
 \int_{\mathbb R^\times}W(a)\overline{W'(a)}\frac{da}{|a|}.
\end{equation}
\begin{lemma}[The \(K_0\) line norm]
\label{v4-arith:kp:lem:Whittaker-is-Fock}
For this precise coefficient normalization,
\(H_{\mathrm{Whit},\infty}(W_0,W_0)=\pi/4\).
\end{lemma}
\begin{proof}
The classical integral representation
\(K_0(x)=\int_0^\infty e^{-x\cosh u}\,du\), \(x>0\)
(NIST DLMF (10.32.9)), and Tonelli give
\[
 \int_0^\infty K_0(x)^2\,dx
 =\int_0^\infty\!\int_0^\infty
       \frac{du\,dv}{\cosh u+\cosh v}.
\]
Put \(a=(u+v)/2\), \(b=(u-v)/2\). The Jacobian is two,
\(a\ge0\), \(|b|\le a\), and the denominator is
\(2\cosh a\cosh b\). If \(F(a)=\int_0^a\operatorname{sech}b\,db\),
the integral is \(2\int_0^\infty F'(a)F(a)\,da=F(\infty)^2=\pi^2/4\).
Consequently
\(\|W_0\|^2=2\int_0^\infty K_0(2\pi a)^2\,da=\pi/4\).
\end{proof}
\begin{proposition}[Normalized archimedean vacuum-line comparison]
\label{v4-arith:kp:prop:Pet-is-Segal-infty}
Let \(\Omega_\infty\) be the unit Segal vacuum. The map
\(U_\infty(a\Omega_\infty)=2aW_0/\sqrt\pi\) is an isometry of
complex lines, with
\begin{equation}
\label{v4-arith:kp:eq:Pet-infty}
 H_{\mathrm{Whit},\infty}(U_\infty v,U_\infty w)
 =H_{\mathrm{Segal}}(v,w)
 \qquad(v,w\in\mathbb C\Omega_\infty).
\end{equation}
\end{proposition}
\begin{proof}
Both sides are \(a\overline b\) for
\(v=a\Omega_\infty,w=b\Omega_\infty\), by the computed norm.
An extension intertwining the full arithmetic or Heisenberg actions
with a Whittaker representation is an additional construction.
Irreducibility statements for different group actions do not produce
such an intertwiner.
\end{proof}

\section{Finite-Place Identification}
\label{v4-arith:kp:finite}

\subsection{Bost--Connes Hilbert Pairing}
\label{v4-arith:kp:finite:BC-Hilbert}

For each prime \(p\), define the weighted occupation model
\(H_{\mathrm{occ},p}=\ell^2(\mathbb N_0,p^{-j})\), with basis
\(\delta_j\) and Hermitian form
\begin{equation}
\label{v4-arith:kp:eq:BC-Hilbert}
 \langle v,w\rangle^{\mathrm{BC}}_p
 =\sum_{j\geq0}p^{-j}v_j\overline{w_j},\qquad
 \langle\delta_j,\delta_k\rangle^{\mathrm{BC}}_p
 =p^{-j}\delta_{jk}.
\end{equation}
The raw shift \(R\delta_j=\delta_{j+1}\) satisfies
\(\|Rv\|=p^{-1/2}\|v\|\); the normalized shift \(p^{1/2}R\)
is an isometry. This definition makes no identification with a
spherical Hecke factor or its arithmetic action. Such an identification
must specify the carrier and intertwining maps separately.

\begin{lemma}[the actual prime GNS carrier]
\label{v4-arith:kp:lem:GNS}
The prime shift Gibbs state at \(\beta=1\) has GNS carrier
\(\operatorname{HS}(\ell^2(\mathbb N_0))\), cyclic vector
\(\rho_p^{1/2}\), and left multiplication representation, where
\(\rho_p e_j=(1-p^{-1})p^{-j}e_j\).
For a specified global critical KMS state, this carrier is
unitarily identified with its prime-generated cyclic subspace.
The weighted occupation pairing
\eqref{v4-arith:kp:eq:BC-Hilbert} is not thereby identified with the
whole GNS representation or an automorphic Plancherel space.
\end{lemma}
\begin{proof}
The matrix-unit construction and the state-preserving cyclic
unitary are Lemma~\ref{v4-arith:kms-gns:lem:GNS-space} and
Theorem~\ref{v4-arith:kms-gns:thm:KMS1-GNS-BC-identification}.
A further identification of pairings requires the maps in
Hypothesis~\ref{v4-arith:kms-gns:hyp:arithmetic-pairing}.
\end{proof}

\subsection{The Koszul Pairing at \texorpdfstring{\(p\)}{p}}
\label{v4-arith:kp:finite:Koszul-p}

An arithmetic local bar--cobar pairing has its own two carriers:
\begin{equation}
\label{v4-arith:kp:eq:Koszul-p}
 B_p^{\mathrm{Ar}}:V_p\otimes V_p^!\longrightarrow\mathbb C.
\end{equation}
Identifying either carrier with a Bost--Connes module or algebra is
part of the arithmetic comparison, not a consequence of the following
weighted Hilbert model.

\begin{proposition}[Bilinear form on the weighted occupation carrier]
\label{v4-arith:kp:prop:Koszul-p-is-BC}
On \(H_{\mathrm{occ}}=\ell^2(\mathbb N_0,p^{-j})\), let \(J\)
be coordinate conjugation and set
\begin{equation}
\label{v4-arith:kp:eq:Koszul-p-is-BC}
 B_{\mathrm{occ}}(v,w)=\sum_{j\ge0}p^{-j}v_jw_j,
 \qquad B_{\mathrm{occ}}(v,Jw)=H_{\mathrm{occ}}(v,w).
\end{equation}
This is a continuous bilinear form. Under
\(D\delta_j=p^{-j/2}E_{jj}\), it equals \(\operatorname{Tr}(DvDw)\).
\end{proposition}
\begin{proof}
Cauchy--Schwarz bounds the absolute series by
\(\|v\|\|w\|\). The diagonal trace is that same series, and
inserting coordinate conjugation gives the Hermitian form.
This constructs the model pairing. It does not identify the
arithmetic bar--cobar dual with vectors \(\mu_{p^j}^*\delta_1\):
for \(j>0\), those vectors are zero in the occupation representation.
Transport of \(B_p^{\mathrm{Ar}}\) requires actual maps on both
carriers and compatibility with their actions.
\end{proof}

\subsection{Bost--Connes vs.~Petersson at \texorpdfstring{\(p\)}{p}}
\label{v4-arith:kp:finite:BC-vs-Pet}

Fix an unramified tempered representation and the spherical
Whittaker normalization \(W(1)=1\). The Kirillov Hermitian form is
\begin{equation}
\label{v4-arith:kp:eq:Pet-p-RS}
 H_{\mathrm{Whit}}(\phi,\phi')
 =\int_{\mathbb Q_p^\times}W_\phi(a)\overline{W_{\phi'}(a)}\,d^\times a.
\end{equation}
Here \(m_p=\operatorname{vol}(\mathbb Z_p^\times)\) fixes the measure.
This local representation-theoretic form is transported to an
arithmetic Petersson realization only by the specified placement maps.

\begin{lemma}[Spherical norm and the compact subgroup volume]
\label{v4-arith:kp:lem:Pet-p-is-Plancherel}
The spherical norm is
\begin{equation}
\label{v4-arith:kp:eq:local-Plancherel}
 C_p=H_{\mathrm{Whit}}(W,W)
 =m_p(1+p^{-1})L_p(1,\operatorname{Ad}\pi_p).
\end{equation}
It is finite, positive, and representation-dependent. In contrast,
the compact subgroup volume for additive-matrix unit measure is
\((1-p^{-1})(1-p^{-2})\).
\end{lemma}
\begin{proof}
Lemma~\ref{v4-arith:bcglue-k5:lem:spherical-norm} proves the norm
formula by the convergent spherical shell series. The compact volume
is computed by reduction modulo \(p\): the invertible matrices form
\((p^2-1)(p^2-p)\) of the \(p^4\) residue classes. These are different
calculations and different measures. An extra factor \(|a|^{-1}\)
in the integral would instead evaluate \(Z_p(0)\); for
\(\alpha=\beta=1\), its shell series is \(m_p\sum(j+1)^2\),
which diverges.
\end{proof}

\begin{proposition}[conditional finite-place pairing comparison]
\label{v4-arith:kp:prop:Pet-p-is-BC-sigma}
Assume the local placement and pairing maps of
Hypothesis~\ref{v4-arith:kms-gns:hyp:arithmetic-pairing}, with common
measure normalization. If they realize the two local forms on the
prime Hilbert--Schmidt carrier by the bilinear trace form \(B_p\)
and \(JX=X^*\), then
\begin{equation}
\label{v4-arith:kp:eq:Pet-p-match}
 \langle\iota_pv,\iota_pw\rangle^{\mathrm{Pet}}_p
 =B_p(c_pv,Jc_pw).
\end{equation}
\end{proposition}
\begin{proof}
On Hilbert--Schmidt operators,
\(B_p(X,JY)=\operatorname{Tr}(XY^*)\), so the stipulated maps
transport the same Hermitian form. The weighted occupation form
has a concrete diagonal realization: for \(m=p^j\), send its
basis vector to \(D_j=p^{-j/2}E_{jj}\). Then
\(\operatorname{Tr}(D_jD_k^*)=\delta_{jk}p^{-j}\) and \(JD_j=D_j\).
The radial-function isometry of
Theorem~\ref{v4-arith:bcglue-k5:thm:K5-absorption-canonical}
uses the chosen multiplicative Haar mass explicitly; the spherical
line is normalized by its full positive norm \(C_p\).
This verifies the normalization on the specified model. The
diagonal subspace is not invariant under left multiplication by
\(V\), and it is not the whole prime GNS representation. Its
arithmetic or automorphic placement remains an assumption.
\end{proof}

\begin{corollary}[Local compatibility under actual pairing maps]
\label{v4-arith:kp:cor:local-compatibility}
At each place, suppose maps on the two arithmetic pairing carriers
identify \(B_p^{\mathrm{Ar}}\) with the bilinear model and transport
the specified conjugation to its model conjugation. Suppose also
that the placed Petersson form is the corresponding Hermitian form.
Then, on that common domain,
\begin{equation}
\label{v4-arith:kp:eq:local-compatibility}
 H_p^{\mathrm{Pet}}(v,w)=B_p^{\mathrm{Ar}}(v,\sigma_p w).
\end{equation}
The archimedean placement is included among these data.
\end{corollary}
\begin{proof}
The model identity is \(H(x,y)=B(x,Jy)\). Apply the supplied
pairing maps and their conjugation compatibility. At finite primes
the weighted diagonal calculation above supplies this identity on
its stated model, without identifying the entire arithmetic carrier.
\end{proof}

\section{Adelic Assembly}
\label{v4-arith:kp:adelic}

\subsection{Tate's Restricted Tensor Product}
\label{v4-arith:kp:adelic:restricted-tensor}

\begin{definition}[Normalized restricted tensor product of pairings]
\label{v4-arith:kp:def:adelic-pairing}
Choose local Hermitian spaces \((H_v,e_v)\) with unit reference
vectors. On elementary tensors equal to these references outside a
finite set, define
\begin{equation}
\label{v4-arith:kp:eq:adelic-restricted-tensor}
 H(\otimes'_v x_v,\otimes'_v y_v)=\prod_v H_v(x_v,y_v).
\end{equation}
Only finitely many factors differ from one. If raw local reference
vectors have nonunit norms, first specify their normalization or
provide a separate global continuation and scalar comparison.
The compact subgroup volume is not the norm of every spherical
vector.
\end{definition}

\begin{lemma}[Tamagawa regularisation of the adelic product]
\label{v4-arith:kp:lem:adelic-convergence}
The ordinary product of the spherical
finite-place volumes for \(GL_{n}\) is not the Tamagawa volume before
regularisation. For \(GL_{n}\) over \(\mathbb{Q}\),
\(\mathrm{vol}(GL_{n}(\mathbb{Z}_{p}))=\prod_{k=1}^{n}(1-p^{-k})\),
so the \(k=1\) factor is the divergent Euler factor
\(\zeta(1)^{-1}\). The factors with \(k\geq2\) converge to
\(\prod_{k=2}^{n}\zeta(k)^{-1}\). The finite non-zero value
\(\tau(GL_{2})=1\) is obtained only after the completed-zeta residue,
the archimedean Haar factor, and the quotient by the rational
character have been inserted.
\end{lemma}

\begin{proof}
Weil 1967 \emph{Basic Number Theory} Ch.~VII \S3 computes
\(\mathrm{vol}([GL_{n}])\) for the Tamagawa normalisation. For
\(GL_{2}\), the Tamagawa number is \(\tau(GL_{2})=1\) by Weil's theorem
(Platonov--Rapinchuk 1994 Ch.~3 \S1.3, Theorem~3.4 for split groups
over \(\mathbb{Q}\)); the computation uses the completed zeta function
and the different in an essential way, not merely a residue.
The full Tamagawa formalism is
classical: see Weil 1967 and the full modern account in
Platonov--Rapinchuk \emph{Algebraic Groups and Number Theory}
(1994) Ch.~3.
\end{proof}

\subsection{Assembly of Petersson and Koszul}
\label{v4-arith:kp:adelic:assembly}

\begin{theorem}[Restricted tensor assembly of normalized local forms]
\label{v4-arith:kp:thm:tate-assembly}
Let \((H_v,e_v)\) be Hilbert spaces with specified unit reference
vectors, conjugations \(J_v e_v=e_v\), and bilinear forms
\(B_v(x,J_vy)=H_v(x,y)\). Their algebraic restricted tensor product
has forms, for elementary tensors,
\begin{align}
\label{v4-arith:kp:eq:Pet-restricted-tensor}
 H(\otimes'_v x_v,\otimes'_v y_v)&=\prod_v H_v(x_v,y_v),\\
\label{v4-arith:kp:eq:Koszul-restricted-tensor}
 B(\otimes'_v x_v,\otimes'_v y_v)&=\prod_v B_v(x_v,y_v).
\end{align}
They extend to the Hilbert completion, and
\(H(x,y)=B(x,Jy)\). Their arithmetic realization requires the actual
local placements, compatible adelic maps and common global scalar
of Hypothesis~\ref{v4-arith:kms-gns:hyp:arithmetic-pairing}.
\end{theorem}
\begin{proof}
Only finitely many factors differ from the reference unit values,
so these products are finite. At each finite stage the Hilbert tensor
form is positive definite and the embeddings obtained by tensoring
with reference vectors are isometries. Their directed union has a
Hilbert completion. Tensor products of the conjugations are
compatible isometries and extend to a conjugation \(J\); define
\(B(x,y)=H(x,Jy)\), which is bounded and agrees with the finite
products. This constructs the restricted tensor model. It does not
identify an arithmetic completed pairing with that model or assign
a value to an unnormalized infinite product at \(s=1\).
\end{proof}

\begin{corollary}[global compatibility via local compatibility]
\label{v4-arith:kp:cor:global-from-local}
Assume the local comparison maps and the completed adelic
realization of Theorem~\ref{v4-arith:kp:thm:tate-assembly}, with
the same specified global normalization for both forms. The global identity
(\ref{v4-arith:kp:eq:main}) follows from the local identities
(\ref{v4-arith:kp:eq:local-compatibility}) by taking the
restricted tensor product. Explicitly,
\begin{equation*}
\langle v,w\rangle^{\mathrm{Pet}}
=\prod_{v}\langle v_{v},w_{v}\rangle^{\mathrm{Pet}}_{v}
=\prod_{v}\langle v_{v},\sigma^{\mathrm{BC}}_{v}(w_{v})\rangle^{\mathrm{Koszul}}_{v}
=\langle v,\sigma^{\mathrm{BC}}(w)\rangle^{\mathrm{Koszul}},
\end{equation*}
using (\ref{v4-arith:kp:eq:Pet-restricted-tensor}),
(\ref{v4-arith:kp:eq:local-compatibility}), and
(\ref{v4-arith:kp:eq:Koszul-restricted-tensor}), with
\(\sigma^{\mathrm{BC}}=\bigotimes_{v}\sigma^{\mathrm{BC}}_{v}\)
the adelic restricted tensor product involution.
\end{corollary}

\section{The Bost--Connes Anti-Involution and its GNS Lift}
\label{v4-arith:kp:BC-involution}

\subsection{Definition and Basic Properties}
\label{v4-arith:kp:BC-involution:def}

The Bost--Connes algebraic anti-involution
\(\tau^{\mathrm{BC}}\) is the adelic extension of its local versions
\(\tau^{\mathrm{BC}}_{p}\). The local anti-involution at a prime \(p\)
is the conjugate-linear anti-automorphism
\begin{equation}
\label{v4-arith:kp:eq:BC-involution-local}
\tau^{\mathrm{BC}}_{p}\colon\mathbb{A}^{\mathrm{BC}}_{p}\to\mathbb{A}^{\mathrm{BC}}_{p},
\quad
\mu_{m}\mapsto\mu_{m}^{*},\quad
\mu_{m}^{*}\mapsto\mu_{m},\quad
e(\gamma)\mapsto e(-\gamma).
\end{equation}
On the prime cyclic GNS space at \(\beta=1\), for \(m=p^k\), the corresponding
antiunitary is not \(\tau^{\mathrm{BC}}_{p}\) itself but
\[
J_{\varphi_{1},p}\pi(\mu_m)\Omega_{\varphi_{1}}
=m^{1/2}\pi(\mu_m^{*})\Omega_{\varphi_{1}},\qquad
J_{\varphi_{1},p}\pi(\mu_m^{*})\Omega_{\varphi_{1}}
=m^{-1/2}\pi(\mu_m)\Omega_{\varphi_{1}} .
\]
The factors are forced by
\(S_{\varphi_{1}}=J_{\varphi_{1}}\Delta_{\varphi_{1}}^{1/2}\). At the
archimedean place, the required antilinear map is coefficient
conjugation in the normalized Segal basis, transported through the
specified arithmetic comparison. It is not the complex-linear identity:
for the unit vacuum \(\Omega\), the Hermitian value
\(H(\Omega,i\Omega)=-i\) whereas the bilinear value
\(B(\Omega,i\Omega)=i\). The relation is
\(H(v,w)=B(v,J_\infty w)\).


\begin{lemma}[properties of the local anti-involution and its lift]
\label{v4-arith:kp:lem:involution-properties}
The local Bost--Connes anti-involution
satisfies:
\begin{itemize}
\item \(\tau^{\mathrm{BC}}_{p}\circ\tau^{\mathrm{BC}}_{p}=\mathrm{id}\)
(involutive; Bost--Connes~\S1.2).
\item \(\tau^{\mathrm{BC}}_{p}\) is \(\mathbb{C}\)-antilinear
(it is the \(*\)-conjugation of the \(C^{*}\)-structure).
\item The GNS lift \(J_{\varphi_{1},p}\) preserves the Bost--Connes
Hilbert pairing:
\(\langle J_{\varphi_{1},p}v,J_{\varphi_{1},p}w\rangle^{\mathrm{BC}}_{p}
=\overline{\langle v,w\rangle^{\mathrm{BC}}_{p}}\).
\end{itemize}
\end{lemma}

\begin{proof}
The first two properties are immediate from the defining formulas
(\ref{v4-arith:kp:eq:BC-involution-local}): involutivity from
\((\mu_{m}^{*})^{*}=\mu_{m}\) and \(e(-(-\gamma))=e(\gamma)\);
\(\mathbb{C}\)-antilinearity from the \(*\)-conjugation structure.
The third is Tomita--Takesaki antiunitarity for the KMS\({}_{1}\)
GNS representation, with the half-density factors displayed above
(Bost--Connes~\S2.2 Thm~25; Connes--Marcolli
arXiv:math/0404128 \S1.8).
\end{proof}

\subsection{Adelic Involution}
\label{v4-arith:kp:BC-involution:adelic}

\begin{definition}[Restricted tensor conjugation]
\label{v4-arith:kp:def:adelic-involution}
For local conjugations \(J_ve_v=e_v\) on the specified normalized
Hilbert models, put \(J=\otimes'_v J_v\) on the algebraic restricted
tensor product. It extends isometrically to its completion.
At infinity use coefficient conjugation on the Segal carrier; at a
prime Gibbs factor use its Hilbert--Schmidt adjoint conjugation.
Transport to the arithmetic carrier requires its specified maps.
\end{definition}
\begin{proposition}[Transport of conjugation to an automorphic domain]
\label{v4-arith:kp:prop:involution-conjugation}
Suppose an isometric placement \(U\) of the completed model into
an automorphic Hilbert space intertwines its local conjugations with
pointwise complex conjugation on a dense restricted tensor domain.
Then this intertwining extends to the completion:
\begin{equation}
\label{v4-arith:kp:eq:involution-conjugation}
 U(Jv)=\overline{U(v)}.
\end{equation}
\end{proposition}
\begin{proof}
Both sides define continuous antilinear maps. Their equality on the
dense domain gives equality everywhere. The local arithmetic
intertwiners are hypotheses here; inversion of Satake parameters
does not construct them from the Bost--Connes algebra involution.
\end{proof}

\section{Spherical-Vector Calculation}
\label{v4-arith:kp:spherical-calc}

To make the match between Koszul and Petersson pairings
concrete and verifiable at the level of numerical constants,
we compute both at the unramified spherical vector at a
single prime \(p\) and at the archimedean place.

\subsection{Spherical Petersson at \texorpdfstring{\(p\)}{p}}
\label{v4-arith:kp:spherical-calc:Pet-p}

\begin{lemma}[Spherical Whittaker value]
\label{v4-arith:kp:lem:spherical-Pet-p}
For the spherical function normalized by \(W(1)=1\),
\begin{equation}
\label{v4-arith:kp:eq:spherical-Pet-p}
 H_{\mathrm{Whit}}(W,W)=C_p
 =m_p(1+p^{-1})L_p(1,\operatorname{Ad}\pi_p).
\end{equation}
\end{lemma}
\begin{proof}
This is Lemma~\ref{v4-arith:kp:lem:Pet-p-is-Plancherel}.
\end{proof}

\subsection{The Normalized Model Line}
\label{v4-arith:kp:spherical-calc:Koszul-p}
\begin{lemma}[Unit reference vectors and bilinear pairing]
\label{v4-arith:kp:lem:spherical-Koszul-p}
The occupation reference vector \(\delta_0\), the Gibbs vector
\(\Omega=\sqrt\rho\), and the normalized spherical vector
\(W/\sqrt{C_p}\) each have norm one. The specified line maps
between them are isometries. On the occupation model,
\begin{equation}
\label{v4-arith:kp:eq:spherical-Koszul-p}
 B_{\mathrm{occ}}(\delta_0,J\delta_0)=1.
\end{equation}
An arithmetic Koszul vacuum has this value only when its comparison
maps carry it to the specified model reference vector.
\end{lemma}
\begin{proof}
The occupation Gram matrix gives the first norm, \(\operatorname{Tr}\rho=1\)
gives the second, and the preceding lemma gives the third.
The displayed bilinear identity follows from the coordinate formula.
\end{proof}

\subsection{Global Reference Normalization}
\label{v4-arith:kp:spherical-calc:global}
\begin{proposition}[Reference value in the restricted tensor model]
\label{v4-arith:kp:prop:global-spherical-match}
For the normalized restricted tensor model of
Theorem~\ref{v4-arith:kp:thm:tate-assembly}, the reference vector
\(e=\otimes'_v e_v\) satisfies
\begin{equation}
\label{v4-arith:kp:eq:global-spherical-match}
 H(e,e)=B(e,Je)=1.
\end{equation}
If its arithmetic realization multiplies both forms by a specified
scalar \(c\), both reference values become \(c\). The scalar and
realization must be constructed for that arithmetic representation.
\end{proposition}
\begin{proof}
Every finite-stage reference vector has norm one and the embeddings
are isometries. The identity therefore holds in the completion;
common scalar multiplication preserves it. A cuspidal Petersson norm
is not determined by the volume of the whole automorphic quotient.
The Tamagawa number consequently does not fix this scalar.
\end{proof}

\section{Analytic Continuation and Poles}
\label{v4-arith:kp:analytic-cont}

\subsection{Pole Structure of Both Pairings}
\label{v4-arith:kp:analytic-cont:poles}

For the specified unramified tempered spherical model, the local
meromorphic function is
\[
 Z_p(s)=m_p\frac{L_p(s,\pi_p\times\widetilde\pi_p)}{\zeta_p(2s)}.
\]
Its rank-four Rankin--Selberg factor is
\(\zeta_p(s)L_p(s,\operatorname{Ad}\pi_p)\); it is not the
rank-two standard factor
\((1-\alpha p^{-s})^{-1}(1-\beta p^{-s})^{-1}\).
A meromorphic arithmetic bar--cobar pairing requires its own family
of carriers, domains and compatible comparison maps.

\begin{proposition}[Pole matching under a meromorphic comparison]
\label{v4-arith:kp:prop:pole-matching}
Suppose two scalar pairings obtained from specified meromorphic
families of vectors extend to meromorphic functions on a connected
parameter domain. If the actual pairing comparison identifies them
on a nonempty open subset, they agree meromorphically, including
their pole orders and residues.
\end{proposition}
\begin{proof}
Their difference is meromorphic and vanishes on an open subset.
The identity theorem, after clearing poles locally, makes it zero
on the connected domain. Equal Laurent expansions give the pole
and residue statement. Neither the local spherical formula nor a
partition function constructs the assumed arithmetic family or its
comparison. In particular no universal assertion that all such
pairings have only the two Tate boundary poles follows.
\end{proof}

\subsection{No Exotic Poles}
\label{v4-arith:kp:analytic-cont:no-exotic}

\begin{remark}[The meromorphic arithmetic comparison]
\label{v4-arith:kp:rem:no-new-residue}
Once the specified meromorphic pairing comparison is constructed,
its common Laurent coefficients determine all residues. Without
that comparison, neither the pole divisor nor its equality with a
Petersson family is supplied by the local Hilbert model. The full
parameter-dependent arithmetic pairing remains the construction
required by Proposition~\ref{v4-arith:kp:prop:pole-matching}.
\end{remark}

\section{The Main Theorem}
\label{v4-arith:kp:main-theorem}

\begin{theorem}[Koszul--Petersson pairing compatibility, domain-restricted]
\label{v4-arith:kp:thm:main-full}
Assume the arithmetic placement, local pairing maps and completed
adelic extension of Hypothesis~\ref{v4-arith:kms-gns:hyp:arithmetic-pairing}
on the stated normalized core.
\emph{Conditional statement.} \emph{(domain-restricted to the normalized
sph-fin cuspidal core and using the Tamagawa regularisation of
Lemma~\ref{v4-arith:kp:lem:adelic-convergence}; see
Remark~\ref{v4-arith:kp:rem:domain-restriction})}. The identity
(\ref{v4-arith:kp:eq:main})
\begin{equation*}
\langle v,w\rangle^{\mathrm{Pet}}
\;=\;
\langle v,\sigma^{\mathrm{BC}}(w)\rangle^{\mathrm{Koszul}}
\end{equation*}
holds for all
\(v,w\in\mathcal{V}^{\mathrm{prim}}_{\mathrm{sph\text{-}fin,cusp}}\),
the normalized sph-fin cuspidal core of \(\mathcal{V}^{\mathrm{prim}}\).
Equivalently, the arithmetic Koszul pairing of the sph-fin
cuspidal core with
\((\mathcal{V}^{\mathrm{prim}})^{!}=\mathbb{A}^{\mathrm{BC}}\)
coincides, after Bost--Connes involution on the dual, with the
adelic Petersson inner product on \(L^{2}([GL_{2}])\) restricted
to the image of the sph-fin cuspidal core. Full-sector
compatibility and KMS\(_{1}\)/GNS identification beyond this core
remain open
(Remark~\ref{v4-arith:kp:rem:domain-restriction}).
\end{theorem}

\begin{proof}
For every place \(v\leq\infty\),
Corollary~\ref{v4-arith:kp:cor:local-compatibility} gives
\(\langle v_{v},w_{v}\rangle^{\mathrm{Pet}}_{v}
=\langle v_{v},\sigma^{\mathrm{BC}}_{v}(w_{v})\rangle^{\mathrm{Koszul}}_{v}\).
Theorem~\ref{v4-arith:kp:thm:tate-assembly} provides the restricted
tensor product assembly of both sides. The adelic involution
\(\sigma^{\mathrm{BC}}=\bigotimes_{v}'\sigma^{\mathrm{BC}}_{v}\)
is well-defined on the restricted tensor product
(Definition~\ref{v4-arith:kp:def:adelic-involution}). Taking the product
over \(v\) of the local identities yields
\begin{equation*}
\langle v,w\rangle^{\mathrm{Pet}}
=\prod_{v}\langle v_{v},w_{v}\rangle^{\mathrm{Pet}}_{v}
=\prod_{v}\langle v_{v},\sigma^{\mathrm{BC}}_{v}(w_{v})\rangle^{\mathrm{Koszul}}_{v}
=\langle v,\sigma^{\mathrm{BC}}(w)\rangle^{\mathrm{Koszul}},
\end{equation*}
as in Corollary~\ref{v4-arith:kp:cor:global-from-local}.

The comparison is first defined on the stipulated normalized
restricted-tensor core. Its extension to the completed arithmetic
carrier is part of Hypothesis~\ref{v4-arith:kms-gns:hyp:arithmetic-pairing}.
The spherical-volume product has the Tamagawa residue normalization
of Lemma~\ref{v4-arith:kp:lem:adelic-convergence}; this does not
assert ordinary absolute convergence at \(s=1\). Where a continued
pairing is used, both its domain and compatibility with the comparison
maps must be supplied. Thus the displayed equality follows from the
specified local identities and their assumed adelic extension.
\end{proof}

\begin{remark}[local forms and their comparison maps]
\label{v4-arith:kp:rem:non-circular}
The spherical calculation fixes the multiplicative Haar measure and
computes the Whittaker norm. The weighted occupation and diagonal
Hilbert--Schmidt models have explicit isometries; at infinity the
normalized vacuum line has the computed \(K_0\)-Whittaker norm.
These are comparisons of the stated models and lines. Neither
Stone--von Neumann uniqueness nor a local functional equation
constructs the arithmetic bar-to-automorphic map. Equality of the
arithmetic pairings requires the placement, conjugation and adelic
extension in Hypothesis~\ref{v4-arith:kms-gns:hyp:arithmetic-pairing}.
\end{remark}

\section{Consequence: P3 on the Sph-Fin Cuspidal Core}
\label{v4-arith:kp:P3-unconditional}

\begin{corollary}[P3 unitarity on the sph-fin cuspidal core]
\label{v4-arith:kp:cor:P3-unconditional}
Assume the arithmetic placement, local pairing maps and completed
adelic extension of Hypothesis~\ref{v4-arith:kms-gns:hyp:arithmetic-pairing}
on the stated normalized core.
\emph{Conditional statement.} \emph{(domain-restricted to the normalized
sph-fin cuspidal core and to the Tamagawa-regularised pairing)}. On the normalized sph-fin cuspidal core
\(\mathcal{V}^{\mathrm{prim}}_{\mathrm{sph\text{-}fin,cusp}}\),
the primary arithmetic Heisenberg is a unitary arithmetic chiral
algebra: it carries a positive-definite Hermitian form
\begin{equation}
\label{v4-arith:kp:eq:Pet-form}
\langle\,\cdot\,,\,\cdot\,\rangle^{\mathrm{Pet}}\colon\mathcal{V}^{\mathrm{prim}}_{\mathrm{sph\text{-}fin,cusp}}\otimes\overline{\mathcal{V}^{\mathrm{prim}}_{\mathrm{sph\text{-}fin,cusp}}}\to\mathbb{C}
\end{equation}
given by the adelic Petersson--Tamagawa inner product, and this
form is compatible with arithmetic Koszul self-duality via
\begin{equation}
\label{v4-arith:kp:eq:Pet-Koszul-compat}
\langle v,w\rangle^{\mathrm{Pet}}
\;=\;\langle v,\sigma^{\mathrm{BC}}(w)\rangle^{\mathrm{Koszul}}.
\end{equation}
Full-sector P3 is not a consequence of the sph-fin cuspidal identity:
the prime KMS\(_1\)/GNS cyclic-subspace comparison is the operator-algebraic statement
of Theorem~\ref{v4-arith:kms-gns:thm:KMS1-GNS-BC-identification}, while the
ramified Rankin--Selberg and Eisenstein terms are separate domain
conditions (Remark~\ref{v4-arith:kp:rem:domain-restriction}).
\end{corollary}

\begin{proof}
Positive-definiteness of \(\langle\,\cdot\,,\,\cdot\,\rangle^{\mathrm{Pet}}\)
was established place by place in the proof of
Theorem~\ref{v4-arith:thm:P3-resolution}: at \(\infty\) by Segal
Fock positivity; at each prime \(p\) by Bost--Connes
KMS\({}_{\beta=1}\) positivity; adelic assembly by Tate 1950
restricted tensor product preserves positive-definiteness.

Compatibility with arithmetic Koszul self-duality is
Theorem~\ref{v4-arith:kp:thm:main-full}, which is established on the
normalised sph-fin cuspidal core. Combining:
\(\mathcal{V}^{\mathrm{prim}}\) carries both a positive-definite
Hermitian form (Petersson) and an arithmetic Koszul self-duality
pairing (arithmetic bar--cobar), and the two are compatible
via Bost--Connes involution.
\end{proof}

\begin{remark}[form restriction for P3]
\label{v4-arith:kp:rem:P3-upgrade}
Theorem~\ref{v4-arith:thm:P3-resolution} is conditional on
Conjecture~\ref{v4-arith:cl:cj:koszul-Petersson-compat}. Its conditional transport is Theorem~\ref{v4-arith:kp:thm:main-full}, a domain-restricted
and Tamagawa-regularised theorem on the normalized sph-fin cuspidal core.
P3 retains the same arithmetic comparison hypotheses: on the sph-fin cuspidal core,
\(\mathcal{V}^{\mathrm{prim}}_{\mathrm{sph\text{-}fin,cusp}}\) is
unitary with Petersson inner product compatible with the Koszul
pairing; full-sector P3 also requires the ramified Rankin--Selberg
and Eisenstein terms recorded in the determinacy comparison residual
Section~\ref{v4-arith:w3-comparison-B:correction-26}.
\end{remark}

\subsection{Impact on the Sufficient RH Implication}
\label{v4-arith:kp:P3-unconditional:rh}

\begin{corollary}[sufficient RH path, four-input conditional]
\label{v4-arith:kp:cor:rh-minimal-updated}
Under the normalization and comparison hypotheses of
Theorem~\ref{v4-arith:kp:thm:main-full}, the sufficient closure path to
RH in the arithmetic formalism
(Theorem~\ref{v4-arith:rh:thm:sufficient-path}) is the four-input
conditional conjunction of
Theorem~\ref{v4-arith:w3-comparison-B:thm:four-input-branch}:
\begin{enumerate}
\item The arithmetic comparison includes the central-defect identity
and detecting Mellin evaluation of
Theorem~\ref{v4-arith:thm:P1-resolution}.
Its co--contra component retains the coefficient requirements of
Remark~\ref{v4-arith:pos:rem:open-sub-items}.
\item Koszul--Petersson compatibility
(Theorem~\ref{v4-arith:kp:thm:main-full}): domain-restricted to the
normalized sph-fin cuspidal core and conditional on the
Tamagawa-regularised assembly; the full-sector extension has separate
ramified Rankin--Selberg and Eisenstein components, while the
KMS\(_1\)/GNS component is the prime cyclic-subspace statement of
Theorem~\ref{v4-arith:kms-gns:thm:KMS1-GNS-BC-identification}.
\item Arithmetic Hilbert--P\'olya H-P\(^{\mathrm{Ar}}\)
(Hypothesis~\ref{v4-arith:rh:hyp:HPAr}): the independent spectral-matching
problem.
\item Arithmetic Virasoro positivity VB\(^{*}\) main direction
(Theorem~\ref{v4-arith:rh:thm:VB-implies-F-RH}): the classical
positive-trace problem, equivalent to RH.
\end{enumerate}
\end{corollary}

\begin{proof}
After the determinacy comparison the sufficient RH path carries four
independent conditional inputs. The preceding theorem reduces the second to
the sph-fin domain and the Tamagawa-regularised assembly; the other
three remain on the
four-input RH branch residual list
(Section~\ref{v4-arith:app:C:residual-list-w2}).
\end{proof}

\begin{remark}[remaining arithmetic comparison problems]
\label{v4-arith:kp:rem:substantial-reestimate}
The conditional pairing theorem leaves the following mathematical
problems:
\begin{description}
\item[Arith.~Positselski] non-formal theorem (Iwasawa-theoretic Positselski).
\item[VB\(^{*}\)] the Li positivity / arithmetic Virasoro bootstrap problem.
\item[\(\mathrm{L}_{\mathrm{ACD}}\)] non-formal theorem (adelic
categorical descent; not on RH path).
\item[F2.c] non-formal theorem (Tate \(\infty\)-categorical restricted
tensor; not on RH path).
\item[F3.BZN\(^{\mathrm{Ar}}\)] non-formal theorem (arithmetic
Ben-Zvi--Nadler; not on RH path).
\end{description}
The pairing conclusion of Theorem~\ref{v4-arith:kp:thm:main-full}
retains its arithmetic placement, local comparison, and adelic
extension hypotheses.
\end{remark}

\section{Pairing Constructions and Arithmetic Comparison}
\label{v4-arith:kp:summary}
The explicit constructions are the Segal bilinear/Hermitian identity,
the normalized \(K_0\) vacuum-line isometry, the weighted occupation
and diagonal Hilbert--Schmidt pairing, and the radial-function
isometry of Theorem~\ref{v4-arith:bcglue-k5:thm:K5-absorption-canonical}.
For \(W(1)=1\), the finite spherical norm is
\(m_p(1+p^{-1})L_p(1,\operatorname{Ad}\pi_p)\); unit normalization
uses this full factor. It is not the compact subgroup volume.

The normalized restricted tensor construction is proved in
Theorem~\ref{v4-arith:kp:thm:tate-assembly}. Its identification with
the actual arithmetic completed pairing, its global scalar, and
its action-compatible automorphic placement remain required data.
The domain-restricted Koszul--Petersson theorem retains these
hypotheses; restricting to a cuspidal sph-fin core does not construct
them. Meromorphic pole matching additionally requires an actual
parameter-dependent comparison on an open domain.

The prime KMS/GNS theorem identifies a prime-generated cyclic
subspace, not the whole arithmetic or automorphic Hilbert space.
Ramified, archimedean full-carrier, and continuous-spectrum
comparisons remain part of the original pairing problem.
The sufficient RH implication still retains all four inputs of
Corollary~\ref{v4-arith:kp:cor:rh-minimal-updated}.
